\subsection{工学設計}
工学設計は、製品やシステムのライフサイクル費用を左右する。ソフトウェア設計でも、どの機能を実現するか、どの品質属性を達成するかによって、保守費用、障害対応費用、運用費用が大きく変わる。
工学の問題は、一つの正解を当てる問題ではないことが多い。多くの場合、複数の解があり、利用可能な資源、費用、標準、技術知識、物理法則、組織の制約などの中で、目的に最も合う実現可能な解を選ぶ。
\subsubsection{工学教育における工学設計}
工学教育の認定では、設計能力が重視される。学生は、複雑で開かれた問題に対して、システム、構成要素、過程を設計できることが求められる。その際、健康・安全、標準、経済、環境、文化、社会的考慮を無視できない。
この考え方はソフトウェアにも当てはまる。たとえば医療、金融、公共、交通、製造のソフトウェアでは、機能が動くだけでは足りない。標準、規制、安全性、社会的影響を設計段階から考える必要がある。
\subsubsection{問題解決活動としての設計}
設計は問題解決活動である。しかし、設計問題は開かれており、曖昧に定義され、複数の解き方があることが多い。このような問題は邪悪な問題とも呼ばれる。ここでの邪悪とは悪意があるという意味ではなく、問題を完全に定義するには、一度解こうとしてみる必要があるという意味である。
ソフトウェア開発では、試作や小さな実装を通じて初めて要求が明確になることがある。したがって、設計は最初から完成した正解を出す活動ではなく、理解と解決を往復しながら進む活動である。
\par\smallskip
\begin{center}
\centering
\IfFileExists{assets/generated_figures/ch18/fig-ch18-03.png}{\includegraphics[width=0.96\linewidth]{assets/generated_figures/ch18/fig-ch18-03.png}}{\fbox{\parbox[c][48mm][c]{0.9\linewidth}{\centering 画像生成待ち\\邪悪な問題としての設計を示す図}}}
\par\smallskip
\refstepcounter{figure}\label{fig:ch18-03}図 \thefigure: 邪悪な問題としての設計を示す図
\end{center}
図 \thefigure は、邪悪な問題としての設計を示す図を視覚的に整理したものである。
\par\smallskip
